\scanpage{123}
\uline{11. Diviseurs relatifs.}

\uline{11.1.} Soit $X$ un $S$-préschéma, $f:X\longrightarrow S$ le morphisme
structural, supposé plat et localement de présentation finie,
d'où le faisceau $\underline R_{X/S}$ de 10.1. Le faisceau quotient
\[
\underline{\operatorname{Div}}_{X/S}\simeq\underline R_{X/S}^{\ast}/\underline O_X^{\ast}
\]
s'appelle le faisceau \uline{des diviseurs de} $X$ \uline{relativement à} $S$, ou
\uline{transversaux à} $f$. Le groupe de ses sections est noté
$\operatorname{Div}(X/S)$ et ses éléments sont appelés \uline{diviseurs de} $X$ \uline{relatifs}
\uline{à} $S$, \uline{où trans transversaux à} $S$. On a les inclusions canoniques
\[
\underline{\operatorname{Div}}_{X/S}\subset\underline{\operatorname{Div}}_X
\]
\[
\operatorname{Div}(X/S)\subset\operatorname{Div}(\unsure{x|X})
\]
donc les diviseurs relatifs sont de cas particuliers de diviseurs.
On a encore pour tout ouvert $U$ de $X$
\[
\underline{\operatorname{Div}}_{X/S}(U)\simeq\operatorname{Div}(U/S).
\]

\uline{11.2.} Si $f\quad R(X/S)^{\ast}$, on note $\operatorname{div}f\quad\operatorname{Div}(X/S)$ la section
image, notations qui est compatible avec celle de 3.1. Les
formules et remarque de 3.2. s'étendent alors au cas relatif
de façon évidente.

\uline{11.3.} On munit $\underline{\operatorname{Div}}_{X/S}$ de la structure d'ordre induit par
$\underline{\operatorname{Div}}_X$, donc $\operatorname{Div}(X/S)$ de celle induit par
$\operatorname{Div}(X)$, notations $\underline{\operatorname{Div}}_{X/S}^{+},\ \operatorname{Div}^{+}(X/S)$.
\mspage{122}
